\section{Introduction}\label{sec:introduction}

An input in row-major order already contains every bit of its transpose.
The difficulty on a tape machine is to deliver those bits in column-major
order while paying for every movement and operation. Routing follows the bits
through their successive locations. Coding instead allows a cell or message
to depend on several input bits, with the desired bits recovered later. A
routing lower bound therefore need not be a computational lower bound.

Brende and Claude Opus 5.5~\cite{FirstPaper2026} formulate the Tape Coding
Conjecture and give sufficient conditions for excluding linear-time
transposition. They also prove head-pass lower bounds and a macroscopic
isoperimetric condition, Theorem~$M^*$, for the work-head trajectories of fast
transposers. Their construction discussion separates complicated head motion
from the unresolved coding computation. This paper develops both sides of that
separation.

Our results fall into two connected groups. The first concerns abstract local
information access, with $P=d^2b$ input bits. Adaptive nonlinear broadcasts confined to original rows
obey separator bounds. Fixed two-coordinate Boolean packets obey exact distance
bounds after arbitrary bijective row recoding. An entropy extension permits
correlated inputs and helper-selected global bijections: bounded stored-row
degree, inverse support and inverse reuse exclude a fixed fractional saving
when input deficit plus helper entropy is $o(P)$. With an $O(db)$ information
budget, the error is explicitly $O(P/\log d)$. A particular two-stage
architecture also requires substantial global helper information.
The adaptive packet extension permits the next query to depend on earlier
answers. For nonlinear banks with local forward dependencies and small
interval-recovery states, leafwise forests give $P-o(P)$ answers; row-sized
state and information budgets give an $O(P/\log\log d)$ error.
Expander shears supply a complementary obstruction: even bounded inverse row
supports and zero forward column radius do not imply small interval-recovery
states. The obstruction holds for every pair of raw and stored half-row sets,
arbitrary joint input/helper laws, and helper-selected within-row bijections.
It charges average branch budgets and the helper's mutual information with
the input; independent public randomness contributes no information charge.
One fixed native machine computes this shear in linear time from raw input
and blank work tapes. Consequently, cheap encoding does not supply the
adaptive theorem's interval-recovery premise.
The second group concerns actual tape execution. Theorem~$M^*$ imposes a positive lower bound on total time
outside locally affine regions of the limiting head paths. This excludes
specified operations and regular recursions. A fully paid prefix supplies the
original cut condition for any payload computation and the limiting-motion conditions
for any linear-time payload.
The subsequent memo~\cite{PaddingMemo2026} removes the added idle tapes and
extends the padding guarantee to a finer entry-cell charge; see
Section~\ref{sec:padding} for its precise range and scope.

\paragraph{What the padding construction shows.}
As a corollary of Theorem~\ref{thm:motion-affine-time} applied to the $M^*$
limits of Theorem~\ref{thm:padding}, the padding's motion lies outside the
excluded class. The padding theorem does not refute the Tape Coding
Conjecture, whose demand and causal-code requirements are additional
conditions. It establishes that the particular geometric inequalities under
consideration can be satisfied without performing the transpose. To obtain a
transposition lower bound, these inequalities must be combined
with an additional restriction, such as one tracking the information required
at the outputs. For a construction, satisfying them leaves encoding and
decoding to be implemented.

\paragraph{Different scopes.}
The finite coding theorems permit general entry width $b$. Their hypotheses
describe abstract broadcasts or row transformations; they are not automatic
properties of arbitrary tape computations. The native asymptotic theorems use
the canonical logarithmic width and the original finite-state machine model.
The motion results and their applications allow negligible entropy deficit
and, where stated, negligible size-dependent control information. The packet
theorems distinguish all-input row recoding from helper-selected global
recoding with locally recoverable original rows. Their restrictions concern
the inverse dependencies and the subsequent packet list; the forward
bijection may mix the whole input. The adaptive forest theorem instead
requires forward column locality and interval recovery; it allows large
coordinate components and only requires termination on the correct inputs.
The expander-shear result separates interval recovery from bounded inverse
support and linear-time encoding. Its native implementation uses canonical
width and even dimensions. Each change of model is made explicit at
the relevant statement.

\paragraph{Relation to the earlier paper.}
The local coding results turn parts of the earlier construction programme into
proved restrictions; they are not improvements to an existing separator or
two-stage theorem there. The affine results strengthen consequences of
Theorem~$M^*$. The paid construction gives a data-free version of the earlier
expanding-motion proposal: it generates and executes the motion in linear
time, before the payload begins. The new shear couples the same expander
family to a fully paid data transformation. Coupling successive transformations
to a transpose decoder with linear total cost remains open.
The scan and recursion results give applications
in the actual machine model. None improves an unrestricted transposition time
bound.

\subsection{Proofs and their interpretation}\label{sec:proof-policy}

\textbf{The proofs of our new results are the Lean proofs.} Every numbered
formal statement names its declarations and source modules. The Lean statements
are authoritative; the English statements are consequences of them, with any
specialization stated explicitly. Each \emph{Outline of the formal proof}
explains the checked argument for human readers. The declaration index in
\cref{app:index} collects these links.

Formal verification establishes propositions about the stated definitions;
model applicability and novelty require separate arguments. Our review checks
the English hypotheses, quantifiers, conclusions and proof outlines against
the formal source. Imported theorems are proved dependencies. The original
paper's four explicitly unformalized background propositions retain that
status when cited here.

The machine definitions and original theorems are imported directly from the
sibling project. The analytic chain additionally uses an unmodified pinned
OpenAI formal source, as specified in \cref{app:verification}. The established
expander mathematics and imported projection theorem are credited in
\cref{sec:related-work}, together with the literature comparisons for our
contributions.
